\chapter{Cohomology, Moduli, and the \texorpdfstring{\(j\)}{j}-Invariant}
\label{ch:cohomology-moduli}

The Jacobian also has a cohomological incarnation.  The exponential sequence of
sheaves
\[
  0\to\Z\to\OO_E\to\OO_E^*\to0
\]
connects integral cycles, holomorphic functions, and line bundles.  Taking
cohomology yields, in degree one, the complex torus
\[
  \Pic^0(E)=H^1(E,\OO_E)/H^1(E,\Z)\cong\C/\Lambda.
\]

\begin{center}
\resizebox{0.78\linewidth}{!}{\classicExpSequenceDiagram}
\end{center}

This chapter's claim is a structural explanation, not a replacement for the
divisor proof.  Divisors make the group visible through zeros and poles; sheaf
cohomology makes the same group visible through line bundles and periods.

\begin{definition}
The \(j\)-invariant of \(E:Y^2Z=4X^3-g_2XZ^2-g_3Z^3\) is
\[
  j(E)=1728\,\frac{g_2^3}{g_2^3-27g_3^2}.
\]
\end{definition}

Scaling a lattice changes \(g_2\) and \(g_3\), but it does not change \(j\).
Thus \(j\) records the isomorphism class of the complex elliptic curve.  Lattice
shapes are parametrized by \(\tau\in\HH\) modulo the modular group
\(\SL_2(\Z)\), and the \(j\)-function identifies the quotient
\(\HH/\SL_2(\Z)\) with the complex line.

\section*{Exercises}
\addcontentsline{toc}{section}{Exercises}

\exerciseintent{Check that the \(j\)-invariant is unchanged by scaling weights.}
\begin{exercise}
If \(g_2\mapsto a^{-4}g_2\) and \(g_3\mapsto a^{-6}g_3\), show that \(j\) is
unchanged.  Compute separately the scaling weights of \(g_2^3\), \(g_3^2\), and
\(\Delta=g_2^3-27g_3^2\).
\end{exercise}

\exerciseintent{Recognize the cohomological Jacobian as a complex vector space modulo a lattice.}
\begin{exercise}
Use the genus-one facts
\[
  H^1(E,\OO_E)\cong\C,\qquad H^1(E,\Z)\cong\Z^2
\]
to write the quotient \(H^1(E,\OO_E)/H^1(E,\Z)\) in the form
\(\C/\Lambda\).  State explicitly what the rank-two lattice \(\Lambda\) is
standing for.
\end{exercise}

\exerciseintent{Compute the extra lattice rotations at the special square and hexagonal shapes.}
\begin{exercise}
For the lattices \(\Z+i\Z\) and \(\Z+e^{2\pi i/3}\Z\), compute which rotations
preserve the lattice: multiplication by \(i\) in the first case and by
\(e^{2\pi i/3}\) in the second.  Record the corresponding rotation orders.
\end{exercise}

\exerciseintent{Separate what cohomology explains from what the chord--tangent proof still needs.}
\begin{exercise}
Which part of the three-face theorem is re-expressed by the quotient
\[
  H^1(E,\OO_E)/H^1(E,\Z)\cong\C/\Lambda?
\]
Write the map or identification it replaces, and list the two extra calculations
still needed for the chord--tangent formula: the line divisor calculation and
the Abel-sum calculation.
\end{exercise}
